\begin{lemma}[Residues and the kernel of the logarithmic connecting map]
\hypertarget{tag-S6_NEW_RESIDUE_TAG}{}\label{pub:I:residue-kernel-rank}
\StatementRecord{pub:I:residue-kernel-rank}{proved}{S6_NEW_RESIDUE_TAG}{\SeriesRef{inv:A1.0}, \SeriesRef{inv:A1.1}, \SeriesRef{inv:A1.2}, \SeriesRef{inv:A2.1}}
For every admissible parameter, the logarithmic connecting map
$\kappa_1:R^1f_*\mathcal E\to\mathcal L\otimes R^2f_*\OO=\OO$
is surjective. Put
$\mathcal K=\ker\kappa_1$ and
$r=\operatorname{rank}H^0(\kappa_1)$. The residue maps give
\[
f_*\Omega_X^1=\Omega_B^1,\qquad
R^1f_*\Omega_X^1=\mathcal K,\qquad
R^2f_*\Omega_X^1=R^2f_*\mathcal E.
\]
Thus the first and third bundles are isomorphic to $\OO(-2)$ and
$\OO\oplus\OO(-1)$, respectively; the latter splitting is \cref{inv:A1.2}.
Moreover,
\[
h^0(\mathcal K)=2-r,\qquad h^1(\mathcal K)=1-r,\qquad
h^{1,2}=2-r,\qquad h^{1,1}=3-r.
\]
\end{lemma}
\begin{proof}
\begin{enumerate}[label=\arabic*.,leftmargin=*]
\item Let $d_{ij}(s)$ be the cover Kodaira--Spencer coefficients in the
eigenframes of \cref{inv:A2.1} after removing $ds$.
The invariant coarse matrix and its cup-product row have the forms
\[
[\kappa_0]=\begin{pmatrix}A&t_jB\\C&D_0\end{pmatrix},\qquad
[\kappa_1]=(-C,A,-D_0,B),
\]
in the ordered basis
$e_0b_0,s^ae_0b_1,s^{m-a}e_1b_0,e_1b_1$. Specifically
$A=sd_{00}/m$, $B=s^{1-a}d_{01}/m$,
$C=s^{1-a}d_{10}/m$, $D_0=s^{m+1-2a}d_{11}/m$.
These are holomorphic invariant functions; the row follows by
wedging the connecting class with $b_0,b_1$.
If the row had no unit at a finite point, every entry would vanish
modulo $t_j$ and $\det\kappa_0$ would vanish to order at least two,
contradicting \cref{inv:A2.1}. Away from the finite points, including
the cusp, $\kappa_0$ is invertible, which gives a unit in the
cup-product row. Thus $\kappa_1$ is surjective.

\item By \cref{inv:A2.1},
$\coker\kappa_0=\C_{p_1}\oplus\C_{p_2}$.
The logarithmic conormal sequence and \cref{inv:A1.2} give
\[
f_*\mathcal A=\OO(1),\quad
0\to T\to R^1f_*\mathcal A\to\mathcal K\to0,\quad
R^2f_*\mathcal A=\OO\oplus\OO(-1),\qquad
T=\C_{p_1}\oplus\C_{p_2}.
\]
Let $\nu:\widetilde D=\mathrm{dP}_6\sqcup S_1\sqcup S_2\to D$
be normalization. The local bases in \cref{inv:A1.1} give the exact
residue sequence
\[
0\to\Omega_X^1\to\mathcal A\to\nu_*\OO_{\widetilde D}\to0.
\]
Its three residue direct images are
$\C_{p_0}\oplus\C_{p_1}\oplus\C_{p_2}$,
$\C_{p_1}\oplus\C_{p_2}$, and zero: the normalization of $W$ is
rational, and a bielliptic surface has
$(h^0,h^1,h^2)(\OO)=(1,1,0)$.
In degree zero the map from $\OO(1)$ is local evaluation with the
nonzero multipliers $1,3,4$. Its kernel is $\OO(-2)$.
Under the base logarithmic inclusion this kernel is $\Omega_B^1$,
and its map to $f_*\Omega_X^1$ is pullback.

\item At $p_j$, in the invariant frame $(b_0,s^ab_1)$ of $G$, residue
of $(dt_j/t_j)(x b_0+y s^ab_1)$ in $H^1(S_j,\OO)$ is
$m_jx[b_0]$. The residue factors through the base logarithmic line
on the covering chart. Naturality of the connecting map gives
$\operatorname{res}_1\kappa_0=0$.
The determinant of $\kappa_0$ has a simple zero, so its reduction
has rank one, and its image is exactly the one-dimensional kernel
$\C s^ab_1$ of this residue. Hence
$\coker\kappa_0\to H^1(S_j,\OO)$ is an isomorphism.
Both torsion summands cancel, giving the stated canonical
identifications of $R^1f_*\Omega_X^1$ and $R^2f_*\Omega_X^1$.

\item The long exact cohomology sequence of
\[
0\to\mathcal K\to\OO^2\oplus\OO(-1)^2
 \xrightarrow{\kappa_1}\OO\to0
\]
gives $h^0(\mathcal K)=2-r$ and $h^1(\mathcal K)=1-r$.
Leray on the curve $B$ and the preceding direct images then give
\[
h^{1,2}=h^1(\mathcal K)+h^0(\OO\oplus\OO(-1))=2-r,
\qquad
h^{1,1}=h^1(\OO(-2))+h^0(\mathcal K)=3-r.
\]
\end{enumerate}
\end{proof}

\begin{proposition}[The global rank and the decisive scalar]\hypertarget{tag-01RX}{}\label{inv:A2.4}\label{wb:thm:h12-scalar}
\StatementRecord{inv:A2.4}{proved}{01RX}{\SeriesRef{pub:I:residue-kernel-rank}, \SeriesRef{inv:A2.1}, \SeriesRef{inv:A2.2}}

The map
$\kappa_1:R^1f_*\mathcal E\to\mathcal L\otimes R^2f_*\OO=\OO$
is surjective, and $H^0(\kappa_1):\C^2\to\C$ has rank one.
Let $h=a\,d\zeta_1+d\zeta_2$ be the unique global section of
$f_*\mathcal E$ with $h(\xi)=1$.  In an invariant target
frame normalized as below,
\[
\kappa_1(h\gamma)(p_2)=-\tfrac14\mu_s(0)=-\tfrac i4 C_4.
\]
The rank-one condition is equivalent to
$a\tau_z+\mu_z\not\equiv0$.  The same direct-image exact sequence expresses the two Hodge numbers
in terms of the global rank $r$ as
$h^{1,2}=2-r$ and $h^{1,1}=3-r$; here $r=1$.
\end{proposition}


\begin{proof}
\begin{enumerate}[label=\arabic*.,leftmargin=*]
\item Sheaf surjectivity is proved in
\cref{pub:I:residue-kernel-rank}. The lift $h$ exists and is unique because $f_*\mathcal E$ is an
extension of $\OO$ by $\OO(-1)$.  Its invariance under $g_2$ gives
$a(g_2z)=\tau a+\mu$, hence $a(z_2)=1/2$.
In the quotient frame $(\gamma,u)$ of $G$ the connecting map is
\[
\kappa_0(h)=\frac{s}{4}
\bigl((6a\mu_s+\beta_s)\gamma+(a\tau_s+\mu_s)u\bigr)\ell.
\]
Choose the variable eigenframe $b_1$ by Reynolds averaging so that
$\gamma\wedge b_1=v(s)\gamma\wedge u$, $v(0)=1$.
Then $s\,v(s)(\gamma\wedge u)\ell$ is the invariant
target frame; the uncorrected expression with $v=1$ need not be
invariant away from zero.  Cup product with $\gamma$ gives at the
central point $-(a\tau_s+\mu_s)/4=-\mu_s/4$, proving rank one.

\item To check the stated equivalence, there is a second global section
\[
q_E=\sigma_1\otimes u+6\sigma_2\otimes\gamma
\quad\text{of }R^1f_*\mathcal E.
\]
It is the negative of the relative $(1,1)$ image of the invariant
flat class $\eta=u\wedge w+6\gamma\wedge\delta$.
It extends at the cusp because the relative logarithmic forms and
$\gamma,u$ extend there.  At a finite point it is an invariant
holomorphic tensor on the smooth cover, hence belongs to the four
explicit invariant generators of \cref{pub:I:residue-kernel-rank}.  Its cup-product row is zero,
since the two terms are $6\mu_z$ and $-6\mu_z$.
It is independent of $h\gamma$ on every regular fibre (its
$\sigma_1\otimes u$ coefficient is nonzero), so these two sections
form a basis of the two-dimensional global section space.
Consequently the global map has rank one precisely when its value
on $h\gamma$, equivalently $a\tau_z+\mu_z$, is not identically
zero. The Hodge-number formulas follow from
\cref{pub:I:residue-kernel-rank}, and the nonzero scalar above gives
$r=1$.
\end{enumerate}
\end{proof}

\subsection{Absolute cotangent direct images}

\begin{proposition}[All absolute cotangent direct images]\hypertarget{tag-01RT}{}\label{inv:A1.3}\label{wb:prop:absolute-cotangent-direct-images}
\StatementRecord{inv:A1.3}{proved}{01RT}{\SeriesRef{pub:I:residue-kernel-rank}, \SeriesRef{inv:A2.4}}

For every admissible $c$, the six direct images are
\begin{equation}\hypertarget{tag-01T2}{}\label{inv:eq:A-absolute-table}
\begin{array}{c|cc}
q&R^qf_*\Omega_X^1&R^qf_*\Omega_X^2\\\hline
0&\OO(-2)&\OO(-2)\oplus\OO(-1)\\
1&\OO\oplus\OO(-1)^2&\OO(-2)\oplus\OO(-1)^2\\
2&\OO\oplus\OO(-1)&\OO.
\end{array}
\end{equation}
In particular every listed sheaf is torsion-free.
\end{proposition}


\begin{proof}
\begin{enumerate}[label=\arabic*.,leftmargin=*]
\item By \cref{pub:I:residue-kernel-rank}, the first-column direct
images are $\OO(-2)$, $\mathcal K=\ker\kappa_1$, and
$\OO\oplus\OO(-1)$. By \cref{inv:A2.4}, the global map
$H^0(\kappa_1)$ has rank one. Hence the bundle map
$\OO^2\oplus\OO(-1)^2\to\OO$ has a nonzero constant component
on a trivial summand. Eliminating all other components by bundle
automorphisms gives $\mathcal K=\OO\oplus\OO(-1)^2$.

\item Finally $\Omega_X^2=(\Omega_X^1)^\vee\otimes K_X$. Relative
duality, with the locally free direct images now established, gives
\[
R^qf_*\Omega_X^2
=\bigl(R^{2-q}f_*\Omega_X^1\bigr)^\vee\otimes K_B.
\]
Substitution of $K_B=\OO(-2)$ gives every entry of the second column.
\end{enumerate}
\end{proof}
